\section*{ОБОЗНАЧЕНИЯ И СОКРАЩЕНИЯ}
\addcontentsline{toc}{section}{ОБОЗНАЧЕНИЯ И СОКРАЩЕНИЯ}

БД~--~база данных;

ИП~--~интернет-провайдер;

ИС~--~информационная система;

ИТ~--~информационные технологии;

ОС~--~операционная система;

ПО~--~программное обеспечение;

ПОП~--~подсистема обработки платежей;

РП~--~рабочий проект;

СМО~--~система массового обслуживания;

СУБД~--~система управления базами данных;

ТЗ~--~техническое задание;

ТП~--~технический проект;

ЭВМ~--~электронно-вычислительная машина;

API (Application Programming Interface)~--~программный интерфейс приложения, набор соглашений и средств взаимодействия одной программной системы с другой;

ARP (Address Resolution Protocol)~--~протокол разрешения адресов, используемый для определения MAC-адреса по известному IP-адресу в локальной сети;

CRM (Customer Relationship Management)~--~система управления взаимоотношениями с клиентами;

CSRF (Cross-Site Request Forgery)~--~межсайтовая подделка запроса, разновидность атаки на веб-приложения;

CSS (Cascading Style Sheets)~--~каскадные таблицы стилей, язык описания внешнего вида веб-страниц;

DNS (Domain Name System)~--~система доменных имен;

DTO (Data Transfer Object)~--~объект передачи данных;

ER (Entity-Relationship)~--~модель <<сущность~--~связь>>, графический способ описания структуры данных;

FTTH (Fiber to the Home)~--~технология подведения оптического волокна до квартиры абонента;

FUP (Fair Usage Policy)~--~политика справедливого использования трафика;

HTML (HyperText Markup Language)~--~язык гипертекстовой разметки веб-страниц;

HTTP (HyperText Transfer Protocol)~--~протокол прикладного уровня для передачи гипертекстовых документов;

HTTPS (HyperText Transfer Protocol Secure)~--~расширение протокола HTTP с шифрованием передаваемых данных средствами TLS;

HTB (Hierarchical Token Bucket)~--~алгоритм иерархического контроля полосы пропускания, применяемый в простых очередях RouterOS;

ISP (Internet Service Provider)~--~поставщик услуг доступа к сети Интернет;

JSON (JavaScript Object Notation)~--~текстовый формат обмена структурированными данными;

MAC (Media Access Control)~--~уникальный аппаратный адрес сетевого интерфейса;

MVC (Model-View-Controller)~--~архитектурный шаблон <<модель~--~представление~--~контроллер>>;

NAT (Network Address Translation)~--~трансляция сетевых адресов;

ORM (Object-Relational Mapping)~--~объектно-реляционное отображение, технология связывания объектов программы со строками реляционной базы данных;

PPPoE (Point-to-Point Protocol over Ethernet)~--~туннельный протокол для установления соединений по схеме <<точка~--~точка>> поверх Ethernet;

QoS (Quality of Service)~--~качество обслуживания, набор технологий приоритизации сетевого трафика;

RBAC (Role-Based Access Control)~--~ролевая модель разграничения доступа;

REST (Representational State Transfer)~--~архитектурный стиль построения распределенных систем, основанный на стандартных операциях протокола HTTP над ресурсами;

RouterOS~--~операционная система маршрутизаторов компании MikroTik;

SLA (Service Level Agreement)~--~соглашение об уровне обслуживания;

SPA (Single Page Application)~--~одностраничное веб-приложение;

SQL (Structured Query Language)~--~язык структурированных запросов к реляционным базам данных;

TCP (Transmission Control Protocol)~--~протокол управления передачей, обеспечивающий гарантированную доставку данных;

UML (Unified Modeling Language)~--~унифицированный язык моделирования, применяемый при разработке программного обеспечения;

URL (Uniform Resource Locator)~--~единый указатель ресурса в сети Интернет;

UTF-8 (Unicode Transformation Format, 8-bit)~--~кодировка символов Юникода переменной длины, обратно совместимая с ASCII;

VLAN (Virtual Local Area Network)~--~виртуальная локальная сеть;

VoIP (Voice over IP)~--~передача голоса по сетям с пакетной коммутацией;

VPN (Virtual Private Network)~--~виртуальная частная сеть;

WebSocket~--~протокол двунаправленного обмена сообщениями поверх TCP, применяемый в браузерных приложениях;

XSS (Cross-Site Scripting)~--~межсайтовый скриптинг, разновидность атаки на веб-приложения через внедрение исполняемого кода в страницы, отображаемые другим пользователям.
